\documentclass[11pt]{article}

% -------------------------------------------------
% Packages
% -------------------------------------------------
\usepackage{amsmath,amssymb,amsthm,mathtools}
\usepackage[T1]{fontenc}
\usepackage{lmodern}
\usepackage{microtype}
\usepackage{enumitem}
\usepackage{hyperref}
\usepackage[nameinlink]{cleveref}
\usepackage{geometry}
\geometry{margin=1in}

% -------------------------------------------------
% Theorem styles
% -------------------------------------------------
\newtheorem{theorem}{Theorem}[section]
\newtheorem{proposition}[theorem]{Proposition}
\newtheorem{corollary}[theorem]{Corollary}
\theoremstyle{definition}
\newtheorem{definition}[theorem]{Definition}
\newtheorem{standingrule}[theorem]{Standing Rule}
\theoremstyle{remark}
\newtheorem{remark}[theorem]{Remark}
\newtheorem{scholium}[theorem]{Scholium}

% -------------------------------------------------
% Status tags
% -------------------------------------------------
\newcommand{\U}{\textnormal{[U]}}
\newcommand{\C}{\textnormal{[C]}}
\newcommand{\D}{\textnormal{[D]}}
\newcommand{\R}{\textnormal{[R]}}

% -------------------------------------------------
% Shortcuts
% -------------------------------------------------
\newcommand{\Pinf}{P_\infty}
\newcommand{\Oinf}{O_\infty}
\newcommand{\Xinf}{X_\infty}
\newcommand{\Sigmaedge}{\Sigma_{p,q}}
\newcommand{\Hedge}{H^{\mathrm{edge}}_{p,q}}
\newcommand{\scopepq}{s_{p,q}}
\newcommand{\WOR}{\mathrm{WOR}}
\newcommand{\POT}{\mathrm{POT}}

% -------------------------------------------------
% Title
% -------------------------------------------------
\title{\textbf{NFC\_WOR\_Branch}\\
Witness-Overlap Recoverability, Promoted Visible Exhaustion, and Branch-Legitimacy Control}
\author{}
\date{}

\begin{document}
\maketitle

\begin{center}
\textit{Elements of Nested Fibrational Cosmology --- Derived Branch}
\end{center}

\tableofcontents

% =================================================
\section*{Front Block}
\addcontentsline{toc}{section}{Front Block}

\subsection*{Imported spine results}
\begin{enumerate}[label=(\arabic*)]
    \item \textbf{Book I:} Observational Quotient Theorem; finite stabilization; defect-bookkeeping discipline.
    \item \textbf{Book III:} persistent observables, transfer maps, transport obstruction, and observational-transfer objects.
    \item \textbf{Book IV:} Universal Source Theorem and the source-side completion ladder ORA--UC.
    \item \textbf{Book V:} Branch Projection Theorem; endpoint visibility discipline; hidden-channel exclusion; Certified Branch Legitimacy.
    \item \textbf{Book VII:} Scoped Certificate Rule; Promotion Requires Transfer Theorem; named failure frontier discipline; canonical closure schema.
\end{enumerate}

\subsection*{Branch-specific data}
\begin{enumerate}[label=(\arabic*)]
    \item A declared promoted WOR object class descended from the universal source through the declared WOR realization machinery.
    \item A declared promoted visible package
    \[
    \Xinf=(\Pinf,\Oinf).
    \]
    \item Persistent realized edge data
    \[
    (\Sigmaedge,\Hedge,\scopepq),
    \qquad
    \scopepq\in\{\mathrm{full},\mathrm{partial},\mathrm{null}\}.
    \]
    \item Promoted equivalence notions: branch-visible, patch-visible, and persistent-edge-visible equivalence.
    \item A declared WOR closure ledger:
    \[
    \text{realization}
    \Rightarrow
    \text{support recovery}
    \Rightarrow
    \text{scope-dynamics control}
    \Rightarrow
    \text{visible exhaustion}
    \Rightarrow
    \text{promoted recoverability}
    \Rightarrow
    \text{collapse/confluence}
    \Rightarrow
    \text{closure-grade status}.
    \]
    \item A named failure frontier consisting of the smallest unresolved promoted burdens.
\end{enumerate}

\subsection*{Endpoint target}
A lawful promoted NFC descendant package in which promoted comparison status, declared visible support, and genuine promoted non-equivalence are governed entirely by declared patch-visible and persistent-edge-visible data, with no hidden-channel supplementation.

\subsection*{Current status}
\[
\boxed{
\text{Lawful promoted branch; theorem-bearing at conditional force; closure-directed; not yet CERT-CLOSE.}
}
\]

\subsection*{Present active frontier}
\[
\boxed{\text{Primary active frontier: UC1e}_0}
\]

\subsection*{Historical frontier note}
\[
\boxed{
\text{Historical primary frontier: UC1b.2', now rewritten and absorbed as a branch theorem at }[C]\text{ force.}
}
\]

\subsection*{Dependency roles}
\[
\boxed{
\text{UC1b.2' controls promoted scope dynamics;}
\qquad
\text{UC1e}_0\text{ controls promoted visible distinctness.}
}
\]

% =================================================
\section{WOR.0 Purpose}

This branch studies lawful promoted witness-overlap recoverability inside the NFC canon. Its task is to determine whether promoted WOR objects admit a certified visible-support package sufficient to govern comparison status, promoted non-equivalence, and downstream collapse/confluence consequences without hidden-channel supplementation.

The WOR branch is structural, not interpretive. It does not claim source-equivalence with any external comparison program. It does not claim project-collapse identity under different labels. It does not license undeclared target-side enrichment. Its force is exactly the force supported by the declared spine imports, the declared promoted support package, and the proved WOR branch theorems.

% =================================================
\section{Imported Spine Structure}

\begin{proposition}[Formation {\U}]
The WOR branch candidate \(B_{\WOR}\) is constitutionally well-formed: its imported source object, declared promoted WOR arena, declared promoted visible package, declared closure ledger, and named failure frontier are all specified.
\end{proposition}

\begin{proposition}[Interpretive Label Admissibility {\U}]
The label ``witness-overlap recoverability,'' as used in this branch, is admissible at the stated structural scope: it adds no theorem-bearing content beyond the declared promoted visible package, the declared support discipline, and the proved WOR branch theorems.
\end{proposition}

\begin{standingrule}[Branch Non-Smuggling {\D}]
No theorem in this book may gain force from undeclared promoted support, hidden carrier enrichment, interpretive supplementation, or external comparison claims not explicitly stated in the import ledger or proved internally in the branch.
\end{standingrule}

\begin{standingrule}[Promotion Discipline {\D}]
No promoted WOR claim may be cited beyond the scope licensed by its proved visible support, declared dependencies, and discharged closure-ledger rung.
\end{standingrule}

\begin{proposition}[Imported Force Discipline {\U}]
Every theorem of the WOR branch carries exactly the force licensed by its declared spine imports, branch-specific hypotheses, and internal proofs; no stronger reading is lawful.
\end{proposition}

% =================================================
\section{Branch-Specific Arena}

\begin{definition}[Promoted WOR Object {\D}]
A promoted WOR object is a lawful descendant comparison object in the declared WOR target arena whose theorem-bearing promoted structure is organized by the visible package
\[
\Xinf=(\Pinf,\Oinf),
\]
together with persistent realized edge data
\[
(\Sigmaedge,\Hedge,\scopepq).
\]
\end{definition}

\begin{definition}[Promoted Visible Package {\D}]
The promoted visible package of a promoted WOR object is the ordered pair
\[
\Xinf=(\Pinf,\Oinf),
\]
where \(\Pinf\) is the promoted patch-visible carrier and \(\Oinf\) is the promoted thin-overlap / persistent-edge-visible carrier.
\end{definition}

\begin{definition}[Persistent Realized Edge Data {\D}]
For a persistent realized edge \(\{p,q\}\), the associated promoted edge data consist of:
\begin{enumerate}[label=(\arabic*)]
    \item the comparison-signature datum \(\Sigmaedge\),
    \item the edge-local history datum \(\Hedge\),
    \item the scope marker
    \[
    \scopepq\in\{\mathrm{full},\mathrm{partial},\mathrm{null}\}.
    \]
\end{enumerate}
\end{definition}

\begin{definition}[Visible Equivalences {\D}]
The WOR branch uses three visible equivalence notions:
\begin{enumerate}[label=(\arabic*)]
    \item promoted branch-visible equivalence,
    \item patch-visible equivalence,
    \item persistent-edge-visible equivalence.
\end{enumerate}
Each theorem in the branch must specify which visible equivalence governs its conclusion.
\end{definition}

\begin{proposition}[Promoted WOR Arena Well-Formedness {\U}]
The promoted WOR arena is well-defined from the declared promoted visible package and persistent realized edge data.
\end{proposition}

\begin{proposition}[No Third Primitive Promoted Layer {\U}]
At promoted primitive scope, no theorem-bearing layer exists beyond the declared visible pair
\[
(\Pinf,\Oinf).
\]
\end{proposition}

\begin{corollary}[Primitive Promoted Exhaustion {\U}]
Any theorem-bearing promoted distinction in the WOR branch must descend through the declared visible package and its persistent-edge refinements.
\end{corollary}

% =================================================
\section{Lawful Descent and Endpoint Visibility}

\begin{definition}[WOR Endpoint Datum {\D}]
The WOR endpoint datum is the promoted descendant package determined by the declared visible support of the promoted WOR object and the declared realization machinery.
\end{definition}

\begin{definition}[WOR Legitimacy Witness {\D}]
A WOR legitimacy witness is certified evidence that the promoted endpoint datum depends only on:
\begin{enumerate}[label=(\arabic*)]
    \item the universal source,
    \item the declared WOR realization machinery,
    \item the declared promoted visible package,
    \item the declared persistent-edge support data,
\end{enumerate}
with no hidden-channel supplementation.
\end{definition}

\begin{proposition}[WOR Branch Projection {\U}]
The promoted WOR endpoint package is a lawful descendant of the universal source under the declared WOR branch projection.
\end{proposition}

\begin{proposition}[WOR Endpoint Visibility {\U}]
Every theorem-bearing promoted WOR distinction is expressible through declared patch-visible or persistent-edge-visible data.
\end{proposition}

\begin{proposition}[No Hidden Endpoint Carrier {\U}]
No theorem-bearing promoted endpoint feature of the WOR branch may depend on structure outside the declared visible package and persistent-edge support data.
\end{proposition}

\begin{theorem}[WOR Legitimacy Theorem {\U}]
The WOR branch is lawful exactly to the extent that its promoted endpoint package descends from the universal source through declared visible support and declared transfer machinery, with no hidden channel required.
\end{theorem}

% =================================================
\section{Support Sufficiency and Promoted Support Recovery}

\begin{definition}[Promoted Support Package {\D}]
The promoted support package of a promoted WOR object consists of
\[
(\Pinf,\Oinf,\Sigmaedge,\Hedge),
\]
together with the declared visible-equivalence discipline and the declared transfer discipline inherited from the spine.
\end{definition}

\begin{definition}[Promoted Support Sufficiency {\D}]
Promoted support sufficiency holds when every lawful promoted WOR distinction is witnessed through the declared promoted support package and no theorem-bearing promoted force requires undeclared supporting structure.
\end{definition}

\begin{proposition}[Promoted Support Sufficiency {\U}]
At promoted scope, the full declared support is exhausted by
\[
(\Pinf,\Oinf,\Sigmaedge,\Hedge)
\]
together with the declared visible-equivalence and transport discipline.
\end{proposition}

\begin{proof}
By Definition 2.2, the primitive promoted visible package of the branch is
\[
\Xinf=(\Pinf,\Oinf).
\]
By Proposition 2.6, no theorem-bearing primitive promoted layer exists beyond this visible pair. Therefore any further theorem-bearing promoted support, if lawful at all, must arise only through the declared persistent-edge refinements of the visible package rather than through a third primitive promoted carrier.

By Section 3, every theorem-bearing promoted WOR distinction must be branch-visible, and no theorem-bearing promoted endpoint feature may depend on a hidden endpoint carrier outside the declared visible package and persistent-edge support data. Thus the branch has no right to introduce additional theorem-bearing promoted support outside the visible package together with its declared edge-level refinements.

By Definition 2.3, the declared theorem-bearing persistent-edge data are the comparison-signature datum \(\Sigmaedge\), the edge-local history datum \(\Hedge\), and the scope marker \(\scopepq\). The marker \(\scopepq\) records lawful comparison status on the edge; it does not constitute an independent support carrier enlarging the declared promoted support package. Therefore the only declared edge-level support carriers are \(\Sigmaedge\) and \(\Hedge\).

Hence, at the stated promoted scope, the full declared support is exhausted by
\[
(\Pinf,\Oinf,\Sigmaedge,\Hedge)
\]
together with the declared visible-equivalence and transport discipline. \qedhere
\end{proof}

\begin{corollary}[No Undeclared Promoted Support {\U}]
No theorem-bearing promoted support may be imported outside the declared promoted support package.
\end{corollary}

\begin{proposition}[Support-Visible Determinacy {\U}]
Whenever two promoted WOR realizations agree on the declared promoted support package up to the governing visible equivalences, they agree on every branch-visible theorem-bearing promoted distinction licensed at the present rung.
\end{proposition}

% =================================================
\section{Promoted Recoverability Machinery}

\begin{theorem}[UC1a: Patch Recoverability {\C}]
In the declared promoted WOR regime, promoted patch-visible data are recoverable up to patch-visible equivalence from the declared support package.
\end{theorem}

\begin{theorem}[UC1d: Persistent-Edge Fixing {\C}]
In the declared promoted WOR regime, persistent realized edge structure is fixed up to persistent-edge-visible equivalence.
\end{theorem}

\begin{proposition}[UC1b.1: Edge-Support Recovery {\C}]
For every persistent realized edge \(\{p,q\}\) in the declared promoted WOR regime, the data
\[
(\Sigmaedge,\Hedge)
\]
are recoverable up to their native edge-visible equivalences.
\end{proposition}

\begin{proposition}[WOR-UC1b.2a: Scope Is Status-Bearing {\C}]
For a persistent realized edge \(\{p,q\}\), a strict change in the scope marker
\[
\scopepq\in\{\mathrm{full},\mathrm{partial},\mathrm{null}\}
\]
is a strict change in lawful comparison status on that edge.
\end{proposition}

\begin{proof}
By branch declaration, the marker \(\scopepq\) is not decorative metadata. It records which lawful comparison scope is certified on the persistent edge. The three values \(\mathrm{full}\), \(\mathrm{partial}\), and \(\mathrm{null}\) are pairwise distinct status values in the declared WOR regime. Therefore a strict change of value in \(\scopepq\) is precisely a strict change in lawful comparison status on the edge. \qedhere
\end{proof}

\begin{theorem}[UC1b.2': Persistent-Edge Hidden-Channel Status Drift Lemma {\C}]
Let \(\{p,q\}\) be an eventually permanent realized edge in a stabilized directed coherent finite WOR system. Assume that on some cofinal tail:
\begin{enumerate}[label=(\arabic*)]
    \item the full declared edge-level support on that edge is exhausted by
    \[
    (\Sigmaedge,\Hedge),
    \]
    \item both \(\Sigmaedge\) and \(\Hedge\) are stable up to their native edge-visible equivalences, and
    \item the scope marker
    \[
    \scopepq\in\{\mathrm{full},\mathrm{partial},\mathrm{null}\}
    \]
    changes strictly.
\end{enumerate}
If neither \(\Sigmaedge\) nor \(\Hedge\) changes in a branch-visible way, then the edge exhibits hidden-channel status drift. Hence such a strict change in \(\scopepq\) is not admissible.

Therefore:
\[
\boxed{
\text{Any lawful strict change in }\scopepq\text{ on a persistent edge requires a branch-visible change in }
\Sigmaedge\text{ or }\Hedge.
}
\]
\end{theorem}

\begin{proof}
By Proposition 5.4a, a strict change in the scope marker \(\scopepq\) is a strict change in lawful comparison status on the edge. Thus the hypothesis that \(\scopepq\) changes strictly is not merely a change of notation or presentation; it is a change in the branch’s own declared lawful comparison-status datum.

By Proposition 4.3, the full declared support available on the persistent edge at the stated scope is exhausted, at edge level, by
\[
(\Sigmaedge,\Hedge)
\]
together with the governing edge-visible equivalence discipline. Here the scope marker \(\scopepq\) records lawful comparison status and is not counted as an independent support carrier; the relevant support carriers at edge level are \(\Sigmaedge\) and \(\Hedge\).

Since both \(\Sigmaedge\) and \(\Hedge\) are stable up to their native edge-visible equivalences on the cofinal tail, no new declared edge-visible support has appeared there.

Therefore the hypothesized strict change in \(\scopepq\) would amount to a strict change in lawful comparison status without any new declared visible support on the edge.

This is precisely hidden-channel status drift in the branch-local sense: new descendant-relevant lawful force appears without new declared support visible through the licensed edge-support package. Such a move is not admissible. By Book V, hidden-channel supplementation is illegitimate; by Book VII, stronger status cannot be promoted without an explicit transferred support theorem. No such transfer theorem is present here.

Hence the hypothesized strict change in \(\scopepq\) is not lawful unless accompanied by a branch-visible change in \(\Sigmaedge\) or \(\Hedge\). \qedhere
\end{proof}

\begin{corollary}[UC1b.3: Edgewise Well-Definedness {\C}]
Persistent edge comparison status is well-defined from the declared visible edge-support package.
\end{corollary}

\begin{theorem}[UC1c: Selector-Elimination Contradiction {\C}]
Any attempt to enrich promoted WOR comparison structure by undeclared selector-type support yields contradiction with promoted support sufficiency and branch non-smuggling.
\end{theorem}

\begin{theorem}[UC1: Promoted Recoverability Theorem {\C}]
In the declared promoted WOR regime, the promoted WOR package is recoverable from the declared support package up to the governing visible equivalences.
\end{theorem}

% =================================================
\section{Promoted Visible Exhaustion}

\begin{definition}[Genuine Promoted Non-Equivalence {\D}]
Two promoted WOR realizations are genuinely promoted non-equivalent if they fail to agree on at least one theorem-bearing promoted distinction licensed by the declared support discipline and governing visible equivalences.
\end{definition}

\begin{definition}[Promoted Patch-Visible Distinction {\D}]
A promoted patch-visible distinction is a genuine promoted non-equivalence witnessed through the promoted patch carrier \(\Pinf\) up to patch-visible equivalence.
\end{definition}

\begin{definition}[Promoted Persistent-Edge-Visible Distinction {\D}]
A promoted persistent-edge-visible distinction is a genuine promoted non-equivalence witnessed through persistent-edge support data
\[
(\Sigmaedge,\Hedge,\scopepq)
\]
up to persistent-edge-visible equivalence.
\end{definition}

\begin{proposition}[Visible Dichotomy at Promoted Scope {\C}]
Any theorem-bearing promoted distinction in the WOR branch must, if genuine, be witnessed either patch-visibly or persistent-edge-visibly.
\end{proposition}

\begin{theorem}[UC1e$_0$: Promoted Patch-or-Edge Separation Theorem {\C}]
Let \(\Xinf\) and \(Y_\infty\) be two promoted WOR realizations in the declared promoted regime. If \(\Xinf\) and \(Y_\infty\) are genuinely promoted non-equivalent, then that non-equivalence is patch-visible or persistent-edge-visible.

Equivalently:
\[
\boxed{
\text{No genuine promoted non-equivalence can remain invisible to both }\Pinf\text{ and the persistent-edge support package.}
}
\]
\end{theorem}

\begin{proof}
Assume \(\Xinf\) and \(Y_\infty\) are genuinely promoted non-equivalent. By Definition 6.1, they fail to agree on at least one theorem-bearing promoted distinction licensed by the branch at the present force.

Suppose, toward contradiction, that this genuine promoted non-equivalence is neither patch-visible nor persistent-edge-visible. Then the branch is committed to a theorem-bearing promoted difference that is not witnessed through the promoted patch-visible carrier \(\Pinf\) and not witnessed through the persistent-edge support data \((\Sigmaedge,\Hedge,\scopepq)\).

But by Proposition 4.3, the full declared promoted support is exhausted by the visible package and its declared persistent-edge refinements. And by Theorem 5.7, every presently licensed theorem-bearing promoted distinction of the WOR package is recoverable from that support up to the governing visible equivalences. Therefore any theorem-bearing promoted difference between \(\Xinf\) and \(Y_\infty\) must already be carried by that declared support package.

The hypothesized non-equivalence, being invisible to both visible channels, is therefore a theorem-bearing promoted difference unsupported by the declared support package. By Corollary 4.4, no such undeclared promoted support is lawful.

This contradiction shows that the assumed genuine promoted non-equivalence cannot remain invisible to both visible channels. Therefore it must be patch-visible or persistent-edge-visible. \qedhere
\end{proof}

\begin{corollary}[UC1e: Promoted Visible Exhaustion {\C}]
No genuine promoted non-equivalence survives outside the declared promoted visible package and its persistent-edge refinements.
\end{corollary}

\begin{proposition}[No Invisible Promoted Residue {\C}]
There is no theorem-bearing promoted residue class capable of supporting genuine promoted non-equivalence beyond the declared patch-visible and persistent-edge-visible layers.
\end{proposition}

% =================================================
\section{Collapse and Confluence}

\begin{definition}[Promoted Collapse Criterion {\D}]
Two promoted WOR realizations satisfy the promoted collapse criterion when they agree on the full declared promoted support package up to the governing visible equivalences.
\end{definition}

\begin{definition}[Promoted Confluence Condition {\D}]
Two lawful promoted descent paths satisfy the promoted confluence condition when they yield promoted endpoint data agreeing on the full declared promoted support package up to the governing visible equivalences.
\end{definition}

\begin{proposition}[Support Agreement Implies No Genuine Visible Distinction {\C}]
If two promoted WOR realizations agree on the declared promoted support package up to the governing visible equivalences, then no patch-visible or persistent-edge-visible distinction separates them.
\end{proposition}

\begin{theorem}[UC2: Promoted Collapse Criterion {\C}]
Let \(\Xinf\) and \(Y_\infty\) be lawful promoted WOR realizations in the declared promoted regime. If \(\Xinf\) and \(Y_\infty\) agree on the full declared promoted support package up to the governing visible equivalences, then they collapse to the same lawful promoted realization.

Equivalently:
\[
\boxed{
\text{Agreement on full declared promoted support forces promoted identity.}
}
\]
\end{theorem}

\begin{proof}
Assume \(\Xinf\) and \(Y_\infty\) agree on the full declared promoted support package up to the governing visible equivalences. Then, by Proposition 7.3, no patch-visible distinction and no persistent-edge-visible distinction separates them.

Suppose, toward contradiction, that \(\Xinf\) and \(Y_\infty\) do not collapse to the same lawful promoted realization. By Definition 6.1, non-collapse here means the persistence of at least one genuine theorem-bearing promoted non-equivalence.

By Theorem 6.5, any genuine promoted non-equivalence in the declared promoted regime must be patch-visible or persistent-edge-visible. But neither such visible distinction is present, by the support-agreement hypothesis and Proposition 7.3. Therefore no genuine promoted non-equivalence remains.

This contradiction shows that \(\Xinf\) and \(Y_\infty\) must collapse to the same lawful promoted realization. \qedhere
\end{proof}

\begin{theorem}[Promoted Confluence Law {\C}]
If two lawful promoted descent paths satisfy the promoted confluence condition, then they yield equivalent promoted endpoint structure.
\end{theorem}

\begin{corollary}[No Genuine Promoted Multiplicity Without Visible Distinction {\C}]
Path multiplicity without patch-visible or persistent-edge-visible distinction does not count as genuine promoted multiplicity.
\end{corollary}

% =================================================
\section{Closure Ledger}

\begin{definition}[WOR Closure Ledger {\D}]
The WOR closure ledger is the ordered sequence
\[
\text{realization}
\Rightarrow
\text{support recovery}
\Rightarrow
\text{scope-dynamics control}
\Rightarrow
\text{visible exhaustion}
\Rightarrow
\text{promoted recoverability}
\Rightarrow
\text{collapse/confluence}
\Rightarrow
\text{closure-grade status}.
\]
\end{definition}

\begin{definition}[Highest Discharged Rung {\D}]
The highest discharged rung of the WOR closure ledger is the strongest rung for which the governing theorem package has been proved with the force declared by its status tag and dependencies.
\end{definition}

\begin{proposition}[Closure-Ledger Monotonicity {\U}]
The lawful force of the WOR branch is monotone in the closure ledger: discharging a higher rung strengthens the lawful branch claim, while failure to discharge a rung blocks promotion beyond it.
\end{proposition}

\begin{proposition}[WOR Closure Ledger Proposition {\U}]
At every stage of the WOR branch, the present lawful force of the branch is exactly the force licensed by its highest discharged closure-ledger rung and no more.
\end{proposition}

\begin{corollary}[No Ledger-Skipping Promotion {\U}]
No lawful promotion to a stronger WOR branch claim is permitted without discharging all prior required closure-ledger rungs or explicitly proving a transfer theorem bypassing them.
\end{corollary}

% =================================================
\section{Named Failure Frontier}

\begin{definition}[Named Failure Frontier {\D}]
The named failure frontier of the WOR branch is the smallest set of unresolved promoted burdens whose discharge is required before the next lawful closure-ledger promotion can occur.
\end{definition}

\begin{proposition}[Historical Primary Frontier {\C}]
Historically, the primary bottleneck of the WOR line was UC1b.2', the promoted scope-dynamics control theorem.
\end{proposition}

\begin{proposition}[Active Primary Frontier {\C}]
At the present manuscript stage, the active primary frontier of the WOR branch is UC1e$_0$, the promoted visible-distinctness theorem.
\end{proposition}

\begin{proposition}[Dependency Roles of the Frontier Items {\C}]
Within the declared WOR closure ledger:
\begin{enumerate}[label=(\arabic*)]
    \item UC1b.2' controls promoted scope dynamics;
    \item UC1e$_0$ controls promoted visible distinctness.
\end{enumerate}
The former now appears in theorem-bearing manuscript form at conditional force; the latter remains the active frontier burden for further branch strengthening.
\end{proposition}

\begin{proposition}[Named Failure Frontier Proposition {\U}]
The smallest presently unresolved promoted burdens of the WOR branch are exactly the named frontier items declared in this section, and no stronger present-status reading may bypass them.
\end{proposition}

% =================================================
\section{Two-Status Reading / Present Status Proposition}

\begin{definition}[Strengthened Reading {\D}]
The strengthened reading of the WOR branch is the strongest lawful claim that would follow if the currently named frontier items were discharged at the required force.
\end{definition}

\begin{definition}[Present Reading {\D}]
The present reading of the WOR branch is the strongest lawful claim supported by the currently discharged closure-ledger rungs together with the named frontier discipline.
\end{definition}

\begin{proposition}[Present Status Proposition {\U}]
The strongest current lawful claim of the WOR branch is exactly the claim licensed by the discharged WOR closure ledger up to the stated named frontier and no more.
\end{proposition}

\begin{corollary}[No Stronger Present Reading {\U}]
No stronger present reading of the WOR branch is lawful unless the named frontier items are discharged or a separate lawful promotion theorem is proved.
\end{corollary}

\begin{proposition}[Strengthened Reading Under Frontier Discharge {\C}]
If UC1e$_0$ is discharged at the required force, then the lawful WOR branch reading advances from promoted support control and conditional collapse/confluence consequences to the next declared closure-ledger level of fully stabilized visible distinctness.
\end{proposition}

% =================================================
\section{Explicit Non-Claims}

\begin{definition}[Explicit Non-Claim {\D}]
An explicit non-claim of the WOR branch is a statement of content that is not licensed by the discharged WOR closure ledger, the declared branch-specific support package, and the proved WOR branch theorems.
\end{definition}

\begin{proposition}[No Source-Equivalence Claim {\U}]
Nothing in this branch book licenses the claim that NFC and any external comparison program are source-equivalent.
\end{proposition}

\begin{proposition}[No Project-Collapse Claim {\U}]
Nothing in this branch book licenses the claim that NFC and any external comparison program are already one project under different labels.
\end{proposition}

\begin{proposition}[No Hidden Promoted Support Claim {\U}]
Nothing in this branch book licenses the existence of theorem-bearing promoted support beyond the declared promoted support package
\[
(\Pinf,\Oinf,\Sigmaedge,\Hedge)
\]
together with the declared visible-equivalence and transfer discipline.
\end{proposition}

\begin{proposition}[No Invisible Promoted Equivalence Claim {\U}]
Nothing in this branch book licenses promoted equivalence stronger than that supported by the discharged visible-exhaustion and collapse/confluence theorems.
\end{proposition}

\begin{proposition}[No Stronger-Than-Ledger Claim {\U}]
No result stronger than the highest discharged WOR closure-ledger rung is licensed by this branch book.
\end{proposition}

\begin{corollary}[Non-Claims Follow the Ledger {\U}]
Every explicit non-claim in this branch is a direct consequence of the present WOR status proposition and the closure-ledger discipline.
\end{corollary}

% =================================================
\section{Theorem-or-Failure Posture}

\begin{definition}[Theorem-or-Failure Posture {\D}]
The theorem-or-failure posture of the WOR branch is the rule that branch force may increase only by discharge of the next required promoted burden or by an explicit lawful transfer theorem; otherwise the branch remains fixed at its present named frontier.
\end{definition}

\begin{proposition}[Frontier Stability Under Non-Discharge {\U}]
If the current named frontier items are not discharged at the required force, then the lawful status of the WOR branch remains fixed at its present reading.
\end{proposition}

\begin{proposition}[Monotone Strengthening Under Discharge {\U}]
If the next named promoted burden is discharged at the required force, then the lawful WOR branch claim strengthens exactly to the next closure-ledger level determined by that burden’s dependency role.
\end{proposition}

\begin{theorem}[WOR Theorem-or-Failure Law {\U}]
At any stated WOR branch scope, exactly one of two outcomes governs the line:
\begin{enumerate}[label=(\arabic*)]
    \item the next named promoted burden is discharged, thereby strengthening the lawful WOR branch claim to the next declared closure-ledger level; or
    \item the branch remains fixed at its current named frontier, with no lawful promotion to stronger status.
\end{enumerate}
\end{theorem}

\begin{corollary}[No Promotion by Compression, Repetition, or Proximity {\U}]
No promoted WOR claim becomes stronger by compression, repetition, adjacency to stronger claims, or interpretive pressure.
\end{corollary}

% =================================================
\section{Canonical Branch-Paper Sufficiency}

\begin{definition}[Canonical WOR Branch-Paper Sufficiency {\D}]
The WOR branch book is canonically sufficient when it states internally:
\begin{enumerate}[label=(\arabic*)]
    \item its imported spine structure,
    \item its promoted branch-specific arena,
    \item its declared promoted support package,
    \item its closure ledger,
    \item its named failure frontier,
    \item its present-status proposition,
    \item its explicit non-claims, and
    \item the theorem-bearing packages required to identify the present lawful frontier and next honest burden,
\end{enumerate}
so that no retired note is required to understand the branch’s present status.
\end{definition}

\begin{proposition}[Internal Status Sufficiency {\U}]
If the WOR branch book states internally its closure ledger, named frontier, and present-status proposition, then the reader can determine the exact present lawful force of the branch without consulting external notes.
\end{proposition}

\begin{proposition}[Internal Frontier Sufficiency {\U}]
If the WOR branch book states internally its named frontier items together with their dependency roles, then the reader can determine the next honest burden of the branch without consulting external notes.
\end{proposition}

\begin{proposition}[Canonical WOR Branch-Paper Sufficiency Proposition {\U}]
The WOR branch book is canonically sufficient if it contains the internal theorem-bearing material needed to determine:
\begin{enumerate}[label=(\arabic*)]
    \item what the branch claims,
    \item what it does not claim,
    \item what its current lawful frontier is, and
    \item what exact burden must be discharged for lawful promotion.
\end{enumerate}
\end{proposition}

\begin{corollary}[No Retired-Note Dependence for Present Status {\U}]
No retired note is required to determine the present WOR branch status once the canonical sufficiency conditions of this section are satisfied.
\end{corollary}

% =================================================
\section{Conclusion / Final Status Proposition}

\begin{proposition}[Final Status Proposition {\U}]
The WOR branch is a lawful promoted structural descendant branch of NFC. Its present lawful endpoint is promoted visible-support control, promoted recoverability at the force licensed by the discharged ledger, and the stated collapse/confluence consequences; it does not yet license stronger claims beyond the named frontier.
\end{proposition}

\begin{corollary}[Present Endpoint Form {\U}]
The present endpoint of the WOR branch is internal promoted structural control, not external merger, source-equivalence, or unrestricted promoted equivalence.
\end{corollary}

\begin{proposition}[Strengthened Future Status Under Frontier Discharge {\C}]
If the active frontier theorem UC1e$_0$ is discharged at the required force, then the WOR branch advances lawfully to the next declared closure-ledger level and acquires the strengthened reading recorded in Section 10.
\end{proposition}

\begin{scholium}[Boundary of the WOR Branch Book {\R}]
This branch book certifies the lawful promoted structural control of witness-overlap recoverability at the stated scope. It does not certify external source-equivalence, project identity, or any promoted claim stronger than the discharged closure-ledger rung.
\end{scholium}

\end{document}
